\chapter{Gromov's invariant and the volume of a hyperbolic manifold}

\section{Gromov's invariant}

\begin{definition}[Singular chain norm]
\label{thurston-ch06-singular-chain-norm}
\group{gromov-invariant}
\source{thurston:page-0136}
\notready
For a real singular chain $c=\sum_i a_i\sigma_i$ on a space $X$, define
\[
\|c\|=\sum_i|a_i|.
\]
\end{definition}

\begin{definition}[Gromov norm]
\label{thurston-ch06-gromov-norm}
\group{gromov-invariant}
\source{thurston:page-0136}
\notready
For $\alpha\in H_k(X;\mathbb R)$, define
\[
\|\alpha\|=\inf\{\|z\|:z\text{ is a singular cycle representing }\alpha\}.
\]
This is generally only a seminorm.
\end{definition}

\begin{proposition}[Functoriality of the Gromov norm]
\label{thurston-ch06-gromov-norm-functoriality}
\group{gromov-invariant}
\source{thurston:page-0136}
\notready
For a continuous map $f:X\to Y$ and $\alpha,\beta\in H_*(X;\mathbb R)$,
\[
\|\alpha+\beta\|\leq\|\alpha\|+\|\beta\|,\qquad
\|\lambda\alpha\|\leq|\lambda|\|\alpha\|,\qquad
\|f_*\alpha\|\leq\|\alpha\|.
\]
Consequently, for a map $f:M_1\to M_2$ of closed oriented manifolds,
\[
\|[M_1]\|\geq|\deg f|\,\|[M_2]\|.
\]
\end{proposition}

\begin{definition}[Hyperbolic straightening]
\label{thurston-ch06-hyperbolic-straightening}
\group{straightening}
\source{thurston:page-0137}
\notready
For vertices $v_0,\ldots,v_k\in H^n$, the straight simplex is the canonical parametrization of their convex hull obtained in the quadratic-form model by central projection of the affine barycentric simplex. A singular simplex in a complete hyperbolic manifold is straightened by lifting it to $H^n$, replacing the lift by the straight simplex with the same vertices, and projecting back. This operation extends linearly to a chain map
\[
\operatorname{straight}:C_*(M)\to C_*(M)
\]
chain-homotopic to the identity and satisfying $\|\operatorname{straight}(c)\|\leq\|c\|$.
\end{definition}

\begin{proposition}[Bounded volume of hyperbolic simplices]
\label{thurston-ch06-simplex-volume-bound}
\group{simplex-volume}
\source{thurston:page-0137,thurston:page-0138,thurston:page-0139}
\notready
For $k\neq1$, there is a finite supremum $v_k$ for the volumes of straight $k$-simplices in $H^n$. In particular,
\[
v_2=\pi,\qquad v_3=1.0149416\ldots.
\]
It is conjectured that $v_k$ is attained by the regular ideal $k$-simplex, and in all dimensions one has
\[
v_k<\frac{v_{k-1}}{k-1}.
\]
\end{proposition}

\begin{proof}
\sketch
Every simplex is contained in an ideal simplex. The ideal triangle has finite area, equal to $\pi$ by Gauss--Bonnet, and the ideal tetrahedron has a unique maximal-volume shape, the regular one. For an ideal $k$-simplex with one vertex at infinity in the upper half-space model, integrate $t^{-k}$ above its $(k-1)$-face. The resulting identity $(k-1)v(\sigma)=\int h^{-(k-1)}dW^{k-1}$ is strictly bounded by the volume of the base simplex, giving the inductive estimate.
\end{proof}

\begin{corollary}[Hyperbolic lower bound]
\label{thurston-ch06-hyperbolic-norm-lower-bound}
\group{gromov-invariant}
\uses{thurston-ch06-simplex-volume-bound}
\source{thurston:page-0139}
\notready
Every closed oriented hyperbolic $n$-manifold with $n>1$ satisfies
\[
\|[M]\|\geq\frac{v(M)}{v_n}.
\]
More generally, if $0\neq\alpha\in H_k(M;\mathbb R)$ with $k\neq1$, and a multiple $\lambda\alpha$ is represented by a $k$-form of pointwise norm at most $1$, then $\|\alpha\|\geq\lambda/v_k$.
\end{corollary}

\begin{definition}[Measure chains and their norm]
\label{thurston-ch06-measure-chains}
\group{measure-chains}
\source{thurston:page-0139,thurston:page-0140}
\notready
For a manifold $M$, a measure $k$-chain is a compactly supported finite signed Borel measure on $C^1(\Delta^k,M)$; let $\mathcal C_k(M)$ denote these chains, with boundary induced by restriction to faces. Its norm is total variation. The resulting homology is the usual homology, and integration of differential forms over measure chains is defined. The measure-chain Gromov norm of a class is the infimum of the norms of representing measure cycles.
\end{definition}

\begin{definition}[Smearing]
\label{thurston-ch06-smearing}
\group{smearing}
\source{thurston:page-0140,thurston:page-0141}
\notready
Let $M=H^n/\Gamma$ be a closed hyperbolic manifold. Normalize Haar measure on $\operatorname{Isom}_+(H^n)$ so that the isometries carrying a fixed point into a region have measure equal to the region's hyperbolic volume. For a simplex $\sigma$ in $H^n$, push this measure from $\Gamma\backslash\operatorname{Isom}_+(H^n)$ to the space of its projected isometric copies in $M$; denote the resulting measure chain by $\operatorname{smear}_M(\sigma)$. Smearing depends only on the isometry class, extends linearly, and is a chain map.
\end{definition}

\begin{theorem}[Gromov proportionality for hyperbolic manifolds]
\label{thurston-ch06-gromov-proportionality}
\group{smearing}
\uses{thurston-ch06-measure-chains,thurston-ch06-smearing,thurston-ch06-hyperbolic-norm-lower-bound}
\source{thurston:page-0140,thurston:page-0141}
\notready
For every closed oriented hyperbolic $n$-manifold,
\[
\|[M]\|=\frac{v(M)}{v_n}.
\]
\end{theorem}

\begin{proof}
\sketch
The lower bound is the differential-form estimate. For a straight simplex $\sigma$ and its reflected copy $\sigma_-$, the cycle $\frac12\operatorname{smear}_M(\sigma-\sigma_-)$ has norm $v(M)$ and integrates the volume form to $v(M)v(\sigma)$, hence represents $v(\sigma)[M]$. Divide by $v(\sigma)$ and take simplices with volumes approaching $v_n$.
\end{proof}

\begin{corollary}[Degree-volume inequality]
\label{thurston-ch06-degree-volume}
\group{smearing}
\uses{thurston-ch06-gromov-proportionality}
\source{thurston:page-0141}
\notready
If $f:M_1\to M_2$ is a map between closed oriented hyperbolic $n$-manifolds, then
\[
v(M_1)\geq|\deg f|v(M_2).
\]
\end{corollary}

\begin{theorem}[Homogeneous-space proportionality]
\label{thurston-ch06-gx-proportionality}
\group{smearing}
\source{thurston:page-0141,thurston:page-0142}
\notready
Let $G$ act transitively on $X$ with compact isotropy groups, and choose an invariant volume form. There is a constant $C$ such that every closed oriented $(G,X)$-manifold $M$ satisfies
\[
\|[M]\|=Cv(M).
\]
If $\alpha\in H_G^*(X)$ is invariant cohomology and $\alpha_M$ is its induced class on $M$, a norm can be chosen on $H_G^*(X)$ so that
\[
\|\operatorname{PD}(\alpha_M)\|=v(M)\|\alpha\|.
\]
\end{theorem}

\begin{lemma}[Mostow volume rigidity]
\label{thurston-ch06-mostow-volume-rigidity}
\group{mostow}
\uses{thurston-ch06-degree-volume}
\source{thurston:page-0142}
\notready
Homotopy-equivalent closed oriented hyperbolic manifolds have equal volume.
\end{lemma}

\begin{theorem}[Strict Gromov rigidity]
\label{thurston-ch06-strict-gromov-rigidity}
\group{mostow}
\uses{thurston-ch06-degree-volume}
\source{thurston:page-0143}
\notready
Let $f:M_1\to M_2$ be a nonzero-degree map between closed oriented hyperbolic $3$-manifolds. If $v(M_1)=|\deg f|v(M_2)$, then $f$ is homotopic to a local isometry; for degree $\pm1$ it is a homotopy equivalence, and otherwise it is homotopic to a covering map.
\end{theorem}

\begin{theorem}[Exponential convergence of regular simplex volumes]
\label{thurston-ch06-regular-simplex-convergence}
\group{mostow}
\source{thurston:page-0143,thurston:page-0144}
\notready
If $\sigma_E\subset H^3$ is regular with all edge lengths $E$, then $v_3-v(\sigma_E)$ decreases exponentially in $E$.
\end{theorem}

\begin{lemma}[Near-maximal simplex angles]
\label{thurston-ch06-near-maximal-angles}
\group{mostow}
\uses{thurston-ch06-simplex-volume-bound}
\source{thurston:page-0144,thurston:page-0145}
\notready
Any simplex whose volume is sufficiently close to $v_3$ has all dihedral angles close to $60^\circ$. Moreover, for a face angle $\beta$ of such a simplex,
\[
v_3-v(\sigma)\geq C\beta^2
\]
for a uniform constant $C>0$.
\end{lemma}

\begin{lemma}[Almost-everywhere boundary extension]
\label{thurston-ch06-measurable-boundary-extension}
\group{mostow}
\uses{thurston-ch06-regular-simplex-convergence,thurston-ch06-near-maximal-angles}
\source{thurston:page-0145,thurston:page-0146}
\notready
Under the equality hypothesis of the strict Gromov rigidity theorem, a lift of $f$ to $H^3$ converges along almost every ray to a point of $S^2_\infty$. Its measurable boundary extension carries almost every positively oriented regular ideal simplex to another positively oriented regular ideal simplex.
\end{lemma}

\begin{proof}
\sketch
Compare pairs of large regular simplices that agree near a basepoint but differ by lengthening one ray. Exponential volume efficiency and the quadratic angle deficit imply that the image face angles form a summable sequence, forcing almost-everywhere convergence. The equality case then prevents degenerate limiting vertex quadruples. Preservation of almost all equilateral triangles in a horospherical plane, together with ergodicity of the dense triangular translation group, makes the boundary map essentially fractional linear and conjugates the holonomy groups.
\end{proof}

\section{Manifolds with boundary}

\begin{definition}[Relative Gromov norm]
\label{thurston-ch06-relative-gromov-norm}
\group{relative-invariant}
\source{thurston:page-0147}
\notready
For a pair $(M,A)$, define $\mathcal C_k(M,A)=\mathcal C_k(M)/\mathcal C_k(A)$ and measure the norm of a relative chain by the total variation outside $A$. The norm of a relative class is the infimum over relative cycles. For a compact oriented manifold with boundary, its invariant is $\|[M,\partial M]\|$.
\end{definition}

\begin{definition}[Boundary-small relative norm]
\label{thurston-ch06-boundary-small-norm}
\group{relative-invariant}
\source{thurston:page-0147}
\notready
For a compact $3$-manifold with torus boundary, set
\[
\|[M,\partial M]\|_0=\liminf_{a\to0}\{\|z\|:z\text{ is a straight representative of }[M,\partial M],\ \|\partial z\|\leq a\}.
\]
This requires $\|[\partial M]\|=0$, which holds for unions of tori.
\end{definition}

\begin{proposition}[Linear filling bound on the torus]
\label{thurston-ch06-torus-filling-bound}
\group{relative-invariant}
\source{thurston:page-0148}
\notready
There is $K>0$ such that every homologically trivial $2$-cycle $z$ in $T^2$ bounds a chain $c$ with $\|c\|\leq K\|z\|$.
\end{proposition}

\begin{theorem}[Relative Gromov proportionality]
\label{thurston-ch06-relative-gromov-proportionality}
\group{relative-invariant}
\uses{thurston-ch06-relative-gromov-norm,thurston-ch06-boundary-small-norm,thurston-ch06-torus-filling-bound,thurston-ch06-simplex-volume-bound}
\source{thurston:page-0149,thurston:page-0150}
\notready
If $M$ is a complete oriented finite-volume hyperbolic $3$-manifold and $\overline M$ is its compactification by torus boundary, then
\[
\|[\overline M,\partial\overline M]\|_0
=\|[\overline M,\partial\overline M]\|
=\frac{v(M)}{v_3}.
\]
\end{theorem}

\begin{theorem}[Hyperbolic submanifold lower bound]
\label{thurston-ch06-hyperbolic-submanifold-bound}
\group{relative-invariant}
\uses{thurston-ch06-relative-gromov-proportionality}
\source{thurston:page-0150,thurston:page-0151}
\notready
If a closed oriented $3$-manifold $M$ contains an embedded finite-volume complete hyperbolic submanifold $H$ whose compactification has incompressible boundary, then
\[
\|[M]\|\geq\frac{v(H)}{v_3}.
\]
\end{theorem}

\begin{theorem}[Strict volume decrease under cusp filling]
\label{thurston-ch06-cusp-filling-volume-decrease}
\group{relative-invariant}
\uses{thurston-ch06-strict-gromov-rigidity,thurston-ch06-relative-gromov-proportionality}
\source{thurston:page-0151}
\notready
If $M_1$ is complete finite-volume hyperbolic and $M_2\neq M_1$ is obtained by replacing some cusps of $M_1$ by solid tori and is hyperbolic, then
\[
v(M_1)>v(M_2).
\]
\end{theorem}

\section{Ordinals and commensurability}

\begin{corollary}[Well-ordering of hyperbolic volumes]
\label{thurston-ch06-volume-well-order}
\group{volume-order}
\uses{thurston-ch06-relative-gromov-proportionality,thurston-ch06-cusp-filling-volume-decrease}
\source{thurston:page-0152}
\notready
The set of volumes of complete hyperbolic $3$-manifolds is a closed well-ordered subset of $\mathbb R_+$, and each volume is realized by only finitely many hyperbolic manifolds. Its order type is $\omega^\omega$; volumes with $k$ cusps occur at ordinal levels of the form $\alpha\cdot\omega^k$. There is a constant $C>0$ such that every $k$-cusped manifold has volume at least $Ck$.
\end{corollary}

\begin{corollary}[Well-ordering for graph-manifold decompositions]
\label{thurston-ch06-graph-volume-well-order}
\group{volume-order}
\uses{thurston-ch06-hyperbolic-submanifold-bound,thurston-ch06-torus-filling-bound}
\source{thurston:page-0153}
\notready
For connected manifolds obtained from Seifert fiber spaces and finite-volume hyperbolic manifolds by gluing along incompressible tori, the values $v_3\|[M]\|_0$ form a closed well-ordered subset of $\mathbb R_+$ of order type $\omega^\omega$.
\end{corollary}

\begin{definition}[Commensurability]
\label{thurston-ch06-commensurability}
\group{commensurability}
\source{thurston:page-0153}
\notready
Discrete groups $\Gamma_1,\Gamma_2\leq\operatorname{Isom}(H^n)$ are commensurable if, after conjugating $\Gamma_1$, their intersection has finite index in both. Hyperbolic manifolds are commensurable if they possess homeomorphic finite-sheeted covers.
\end{definition}

\begin{proposition}[Commensurable manifolds admit hyperbolic models]
\label{thurston-ch06-commensurable-hyperbolic-model}
\group{commensurability}
\uses{thurston-ch06-commensurability}
\source{thurston:page-0154,thurston:page-0155}
\notready
If a finite-volume hyperbolic manifold $M_1$ is commensurable with $M_2$, then $M_2$ is homotopy equivalent to a complete finite-volume hyperbolic manifold. In particular, commensurable manifolds have volumes with rational ratio.
\end{proposition}

\begin{definition}[Cusp modulus]
\label{thurston-ch06-cusp-modulus}
\group{commensurability}
\source{thurston:page-0155,thurston:page-0156}
\notready
For a finite-volume hyperbolic $3$-manifold, choose a cusp generator $z\mapsto z+1$ and write a second generator as $z\mapsto z+\alpha$. The cusp modulus is defined up to
\[
\alpha\sim\frac{n\alpha+m}{p\alpha+q},\qquad n,m,p,q\in\mathbb Z,\quad nq-mp\neq0.
\]
Equivalent moduli generate the same field, and the moduli are algebraic numbers.
\end{definition}

\begin{proposition}[Algebraic holonomy]
\label{thurston-ch06-algebraic-holonomy}
\group{commensurability}
\uses{thurston-ch06-cusp-modulus}
\source{thurston:page-0156}
\notready
If $\Gamma\leq\operatorname{PSL}(2,\mathbb C)$ is discrete and $H^3/\Gamma$ has finite volume, then $\Gamma$ is conjugate to a group of matrices with algebraic entries.
\end{proposition}

\begin{theorem}[Bass arithmetic alternative]
\label{thurston-ch06-bass-arithmetic-alternative}
\group{commensurability}
\source{thurston:page-0156}
\notready
For a complete finite-volume hyperbolic $3$-manifold $M$, either $\pi_1(M)$ is conjugate into $\operatorname{PSL}(2,\mathcal O)$, where $\mathcal O$ is the ring of algebraic integers, or $M$ contains a closed incompressible surface not homotopic to a cusp. The arithmetic-integrality condition is equivalent to finite generation of the additive matrix group and to integrality of all traces.
\end{theorem}

\begin{conjecture}[Commensurability maxima and volume lattice]
\label{thurston-ch06-commensurability-questions}
\group{commensurability}
\source{thurston:page-0157}
\notready
It is unknown whether every commensurability class of discrete subgroups of $\operatorname{PSL}(2,\mathbb C)$ has finitely many maximal groups up to isomorphism, and whether the volumes in a fixed class form a discrete set of multiples of one positive number.
\end{conjecture}
